\begin{table}[htbp]
\centering
\small
\begin{tabular}{rllrrlrrrr}
\toprule
rule & tuple & class & refl. & flip & $N$ & $n_{\wcc}$ & $D_{\max}$ & $n_1$ & $n_{\wcc}D_{\max}/2^N$ \\
\midrule
51 & \rt{VVVV} & ergodic & 51 & 51 & $6$--$22$ & 1 & 4\,194\,304 & 0 & 1 \\
54 & \rt{IVVV} & ergodic & 54 & 147 & $6$--$22$ & 2 & 4\,194\,303 & 1 & 2 \\
57 & \rt{VIVV} & ergodic & 99 & 99 & $6$--$22$ & 1 & 4\,194\,304 & 0 & 1 \\
99 & \rt{VVIV} & ergodic & 57 & 57 & $6$--$22$ & 1 & 4\,194\,304 & 0 & 1 \\
147 & \rt{VVVI} & ergodic & 147 & 54 & $6$--$22$ & 1 & 4\,194\,304 & 0 & 1 \\
153 & \rt{VIVI} & ergodic & 195 & 102 & $6$--$22$ & 1 & 4\,194\,304 & 0 & 1 \\
195 & \rt{VVII} & ergodic & 153 & 60 & $6$--$22$ & 1 & 4\,194\,304 & 0 & 1 \\
\midrule
60 & \rt{IIVV} & polynomial & 102 & 195 & $6$--$22$ & 23 & 2\,097\,152 & 2 & 11.5 \\
102 & \rt{IVIV} & polynomial & 60 & 153 & $6$--$22$ & 23 & 2\,097\,152 & 2 & 11.5 \\
105 & \rt{VIIV} & polynomial & 105 & 105 & $6$--$22$ & 12 & 1\,352\,078 & 1 & 3.868 \\
150 & \rt{IVVI} & polynomial & 150 & 150 & $6$--$22$ & 12 & 1\,352\,078 & 1 & 3.868 \\
\midrule
108 & \rt{IIIV} & exponential & 108 & 201 & $6$--$22$ & 299\,426 & 17\,711 & 54\,289 & 1264 \\
156 & \rt{IIVI} & exponential & 198 & 198 & $6$--$22$ & 46\,368 & 576 & 12 & 6.368 \\
198 & \rt{IVII} & exponential & 156 & 156 & $6$--$22$ & 46\,368 & 576 & 12 & 6.368 \\
201 & \rt{VIII} & exponential & 201 & 108 & $6$--$22$ & 170\,625 & 46\,368 & 20\,736 & 1886 \\
\midrule
204 & \rt{IIII} & frozen & 204 & 204 & $6$--$22$ & 4\,194\,304 & 1 & 4\,194\,304 & 1 \\
\bottomrule
\end{tabular}
\caption{The unitary census at \rt{obc0}: all sixteen rules, the reflection and spin-flip partner of each, and the three counts at the largest computed $N$. The last column is the hyperbola $n_{\wcc}D_{\max}\ge 2^N$; it equals one exactly for the six ergodic rules and for the identity, the only rules whose sectors all have the same size, and exceeds it by up to three orders of magnitude for the fragmented ones. Reflection preserves the series; the spin flip does not, and it is what puts an ergodic rule opposite a fragmented one in the pairs $(60,195)$ and $(102,153)$.}
\label{tab:r19census}
\end{table}
